\documentclass[10pt]{extarticle}
\usepackage[a4paper,margin=1.5cm]{geometry}
\usepackage[T1]{fontenc}
\usepackage{amsmath,amssymb}
\usepackage{graphicx}
\usepackage[export]{adjustbox}
\usepackage{multicol}
\usepackage{enumitem}
\usepackage{booktabs,array,tabularx}
\setlength{\columnsep}{0.7cm}
\setlength{\parindent}{0pt}
\setlength{\parskip}{3pt}
\pagestyle{empty}
\makeatletter
% figure registry (figdefs.tex is generated by the extract step)
\newcommand{\qfigdef}[4]{\@namedef{qf@file@#1}{#2}\@namedef{qf@w@#1}{#3}\@namedef{qf@wide@#1}{#4}}
\newcommand{\qf@check}[1]{\@ifundefined{qf@file@#1}{\PackageError{score2}{Unknown figure id #1}{}}{}}
\newcommand{\qf@img}[4][]{\includegraphics[width=#2\dimexpr\@nameuse{qf@w@#4}pt\relax,max width=#3,max height=0.45\textheight,#1]{\@nameuse{qf@file@#4}}}
% Full-width questions: finalize wraps every question that contains a wide figure in \qfullbegin...\qfullend,
% so the WHOLE question (stem, figures, options) is set across the page in reading order, then the two columns
% resume. The question's head - its number, stem and first figure (\qfullsplit marks the end of the head) -
% is kept on one page, so a number is never left at a page foot with its figure overleaf; the rest may break
% across pages normally, so no big gaps are left. Without a \qfullsplit, the whole question is kept together
% if it is under 45% of a page. Never write these yourself.
\newif\ifqf@full
\newif\ifqf@boxed
\newsavebox{\qf@box}
\newcommand{\qf@open}{\global\qf@fulltrue\par\global\qf@boxedtrue
  \setbox\qf@box\vbox\bgroup\hsize\textwidth\linewidth\textwidth}
\newcommand{\qf@emit}[1]{\par\egroup\global\qf@boxedfalse
  \ifdim\dimexpr\ht\qf@box+\dp\qf@box\relax<#1\textheight\noindent\usebox{\qf@box}\par\else\unvbox\qf@box\fi}
\newcommand{\qfullsplit}{\ifqf@boxed\qf@emit{0.85}\fi}
\newcommand{\qf@close}{\ifqf@boxed\qf@emit{0.45}\fi\global\qf@fullfalse}
\newcommand{\qfullbegin}{\end{multicols}\qf@open}
\newcommand{\qfullend}{\qf@close\begin{multicols}{2}}
\newcommand{\qfullstart}{\qf@open}
\newcommand{\qfulldone}{\qf@close}
% \qfig[scale]{id}: block figure, centred, directly where it is written (size set by the tool).
% Fallback if a wide figure is not inside a full-width question: break out of the columns in place - no floats.
\newcommand{\qfig}[2][1]{\qf@check{#2}%
  \ifqf@full\qf@blk{#1}{#2}\else
  \ifnum\@nameuse{qf@wide@#2}=1 \expandafter\@firstoftwo\else\expandafter\@secondoftwo\fi
  {\end{multicols}\begin{center}\qf@img{#1}{\textwidth}{#2}\end{center}\begin{multicols}{2}}%
  {\qf@blk{#1}{#2}}\fi}
\newcommand{\qf@blk}[2]{\par\smallskip{\centering\qf@img{#1}{\linewidth}{#2}\par}\smallskip}
% \qfigi[scale]{id}: inline figure (inside option lists / tables / minipages); never breaks columns.
\newcommand{\qfigi}[2][1]{\qf@check{#2}\qf@img[valign=c]{#1}{\linewidth}{#2}}
\makeatother
% structure macros - lint reads these to check completeness and figure placement
\newcommand{\papersec}[1]{\par\bigskip{\centering\Large\bfseries PAPER-#1\par}\medskip}
\newcommand{\partsec}[2]{\par\medskip\noindent{\large\bfseries PART-#1: #2}\par\smallskip}
\newcommand{\sect}[1]{\par\smallskip\noindent\textbf{#1}\par}
\newcommand{\qn}[1]{\par\smallskip\noindent\textbf{#1.}~}
\newcommand{\paraq}[2]{\par\smallskip\noindent\textbf{Paragraph for Questions #1 and #2}\par}
\newenvironment{qoptions}{\begin{itemize}[leftmargin=1.8em,itemsep=1pt,topsep=2pt,label={}]}{\end{itemize}}
\newcommand{\opt}[2]{\item[(#1)] #2}
% figure registry (generated by score2_tools.py finalize)
\qfigdef{p1_phy_s1_q01_1}{TMAS - Test 03 (10 Sep 2023)-img/p1_phy_s1_q01_1.png}{108.2}{0}
\qfigdef{p1_phy_s1_q02_1}{TMAS - Test 03 (10 Sep 2023)-img/p1_phy_s1_q02_1.png}{108.3}{0}
\qfigdef{p1_phy_s1_q04_1}{TMAS - Test 03 (10 Sep 2023)-img/p1_phy_s1_q04_1.png}{108.3}{0}
\qfigdef{p1_phy_s1_q05_1}{TMAS - Test 03 (10 Sep 2023)-img/p1_phy_s1_q05_1.png}{108.2}{0}
\qfigdef{p1_phy_s2_para7_8_1}{TMAS - Test 03 (10 Sep 2023)-img/p1_phy_s2_para7_8_1.png}{108.5}{0}
\qfigdef{p1_phy_s3_q02_1}{TMAS - Test 03 (10 Sep 2023)-img/p1_phy_s3_q02_1.png}{108.0}{0}
\qfigdef{p1_phy_s3_q03_1}{TMAS - Test 03 (10 Sep 2023)-img/p1_phy_s3_q03_1.png}{108.0}{0}
\qfigdef{p1_phy_s3_q04_1}{TMAS - Test 03 (10 Sep 2023)-img/p1_phy_s3_q04_1.png}{108.0}{0}
\qfigdef{p1_phy_s3_q06_1}{TMAS - Test 03 (10 Sep 2023)-img/p1_phy_s3_q06_1.png}{107.8}{0}
\qfigdef{p1_phy_s3_q07_1}{TMAS - Test 03 (10 Sep 2023)-img/p1_phy_s3_q07_1.png}{108.3}{0}
\qfigdef{p2_phy_s1_para1_2_1}{TMAS - Test 03 (10 Sep 2023)-img/p2_phy_s1_para1_2_1.png}{91.1}{0}
\qfigdef{p2_phy_s1_q01_A}{TMAS - Test 03 (10 Sep 2023)-img/p2_phy_s1_q01_A.png}{97.7}{0}
\qfigdef{p2_phy_s1_q01_B}{TMAS - Test 03 (10 Sep 2023)-img/p2_phy_s1_q01_B.png}{97.5}{0}
\qfigdef{p2_phy_s1_q01_C}{TMAS - Test 03 (10 Sep 2023)-img/p2_phy_s1_q01_C.png}{97.7}{0}
\qfigdef{p2_phy_s1_q01_D}{TMAS - Test 03 (10 Sep 2023)-img/p2_phy_s1_q01_D.png}{97.5}{0}
\qfigdef{p2_phy_s1_para3_4_1}{TMAS - Test 03 (10 Sep 2023)-img/p2_phy_s1_para3_4_1.png}{140.7}{0}
\qfigdef{p2_phy_s1_para5_6_1}{TMAS - Test 03 (10 Sep 2023)-img/p2_phy_s1_para5_6_1.png}{65.1}{0}
\qfigdef{p2_phy_s2_para7_8_1}{TMAS - Test 03 (10 Sep 2023)-img/p2_phy_s2_para7_8_1.png}{124.6}{0}
\qfigdef{p2_phy_s2_para9_10_1}{TMAS - Test 03 (10 Sep 2023)-img/p2_phy_s2_para9_10_1.png}{129.8}{0}
\qfigdef{p2_phy_s2_para11_12_1}{TMAS - Test 03 (10 Sep 2023)-img/p2_phy_s2_para11_12_1.png}{158.2}{0}
\qfigdef{p2_phy_s4_q02_A}{TMAS - Test 03 (10 Sep 2023)-img/p2_phy_s4_q02_A.png}{108.5}{0}
\qfigdef{p2_phy_s4_q02_B}{TMAS - Test 03 (10 Sep 2023)-img/p2_phy_s4_q02_B.png}{108.5}{0}
\qfigdef{p2_phy_s4_q02_C}{TMAS - Test 03 (10 Sep 2023)-img/p2_phy_s4_q02_C.png}{97.9}{0}
\qfigdef{p2_phy_s4_q02_D}{TMAS - Test 03 (10 Sep 2023)-img/p2_phy_s4_q02_D.png}{97.9}{0}
\raggedcolumns
\makeatletter\@beginparpenalty=10000\makeatother
\newcolumntype{L}[1]{>{\raggedright\arraybackslash}p{#1}}
\begin{document}
\begin{center}{\Large\bfseries TMAS --- Test 03 (10 Sep 2023)}\end{center}
\begin{multicols}{2}
\raggedcolumns
\papersec{1}
\partsec{1}{PHYSICS}
\sect{SECTION-I(i)}
\qn{1} A bimetallic strip is formed out of two identical strips, one of copper and other of brass. The coefficient of linear expansion of the two metals are $\alpha_c$ and $\alpha_B$ respectively. On heating, the temperature of the strip increases up by $\Delta T$ and the strip bends to form an arc of radius of curvature $R$. Then $R$ is :-
\qfig{p1_phy_s1_q01_1}
\begin{qoptions}
\opt{A}{Proportional to $\Delta T$.}
\opt{B}{Inversely proportional to $\Delta T$.}
\opt{C}{Proportional to $|\alpha_B-\alpha_C|$.}
\opt{D}{Inversely proportional to $|\alpha_B-\alpha_C|$.}
\end{qoptions}
\qn{2} The tension in a stretched string fixed at both ends is changed by $2\%$, the fundamental frequency is found to get changed by $15$ Hz. Select the correct statement(s).
\qfig{p1_phy_s1_q02_1}
\begin{qoptions}
\opt{A}{Wavelength of fundamental frequency does not change.}
\opt{B}{Velocity of propagation of wave changes by $2\%$.}
\opt{C}{Velocity of propagation of wave changes by $1\%$.}
\opt{D}{Original frequency is $1500$ Hz.}
\end{qoptions}
\qn{3} Two pulses travelling on same string are described by equations $y_1=\dfrac{5}{(3x-4t)^2+2}$ and $y_2=\dfrac{-5}{(3x+4t-6)^2+2}$. Mark out correct statements.
\begin{qoptions}
\opt{A}{At $t=\dfrac43$s the two waves cancel each other everywhere.}
\opt{B}{At $t=\dfrac34$s the two waves cancel each other everywhere.}
\opt{C}{At $x=1$ the two waves cancel each other at all times.}
\opt{D}{At all points two waves cancel each other for all values of $t$.}
\end{qoptions}
\qn{4} A wire of mass $M$, length $L$ and cross-section $A$ is shaped into a ring. It is placed on smooth horizontal surface and is rotated about its axis with constant angular speed $\omega$ as shown. If Young's modulus of material of wire is $Y$, then :-
\qfig{p1_phy_s1_q04_1}
\begin{qoptions}
\opt{A}{Tension developed in the wire is $\dfrac{M\omega^2L}{\pi^2}$.}
\opt{B}{Elongation produced in the wire is $\dfrac{M\omega^2L^2}{4\pi^2YA}$.}
\opt{C}{Stress developed in the wire is $\dfrac{M\omega^2L}{4\pi^2A}$.}
\opt{D}{Strain developed in the wire is $\dfrac{M\omega^2L}{\pi^2YA}$.}
\end{qoptions}
\qn{5} In the circuit shown, the switch has been in position 1 for a long time. Now the switch is shifted to position 2. If this instant is taken as $t=0$, then at time $t=\pi/2$ :-
\qfig{p1_phy_s1_q05_1}
\begin{qoptions}
\opt{A}{Current through the inductor is zero.}
\opt{B}{Current through the inductor is $4$ A.}
\opt{C}{Charge on the capacitor is $4$ C.}
\opt{D}{Potential difference across the capacitor is $4$ V.}
\end{qoptions}
\qn{6} Three simple harmonic motions about the same point in the same direction having the same amplitude $a$ and same period are superposed. If the $3^{\mathrm{rd}}$ SHM leads the second by $45^\circ$ and the second leads the first by same phase angle, then :-
\begin{qoptions}
\opt{A}{The resultant amplitude is $(1+\sqrt2)a$.}
\opt{B}{The phase of the resultant motion relative to the first is $90^\circ$.}
\opt{C}{The energy associated with the resulting motion is $(3+2\sqrt2)$ times the energy associated with any single motion.}
\opt{D}{The resulting motion is not simple harmonic.}
\end{qoptions}
\sect{SECTION-I(ii)}
\paraq{7}{8}
Consider AC circuit given below
\qfig{p1_phy_s2_para7_8_1}
\qn{7} $\omega$ for which average power given by source is maximum, would be :-
\begin{qoptions}
\opt{A}{$\dfrac{1}{\sqrt{LC-R^2C^2}}$}
\opt{B}{$\dfrac{1}{\sqrt{LC}}$}
\opt{C}{$\dfrac{1}{\sqrt{\frac LC-R}}$}
\opt{D}{$\dfrac1R\sqrt{\dfrac LC}$}
\end{qoptions}
\qn{8} If reading of $v_2$ and $v_3$ is $30$ volts and $50$ volts respectively, then reading of $v_1$ would be :-
\begin{qoptions}
\opt{A}{$20$ V}
\opt{B}{$20\sqrt2$ V}
\opt{C}{$10\sqrt2$ V}
\opt{D}{$40$ V}
\end{qoptions}
\paraq{9}{10}
A hollow cylinder has length $L$, inner radius $a$ and outer radius $b$. Thermal conductivity is $k$. Inner and outer surfaces of the cylinder are maintained at temperatures $100^\circ$C and $0^\circ$C respectively. If $kL=11$ J/s, $a=0.1$ m, $b=0.3$ m, and $\ell n3=1.1$.
\qn{9} The rate of radial heat (use $\pi=3.14$) flow (in kJ/s) is :-
\begin{qoptions}
\opt{A}{$8.26$}
\opt{B}{$6.28$}
\opt{C}{$5.75$}
\opt{D}{$7.65$}
\end{qoptions}
\qn{10} If the temperature gradient near the inner surface of cylinder is $\dfrac{1000n}{11}\ {}^\circ$C, then value of $n$ is :-
\begin{qoptions}
\opt{A}{$10$}
\opt{B}{$15$}
\opt{C}{$20$}
\opt{D}{$25$}
\end{qoptions}
\sect{SECTION-II}
\qn{1} A certain amount of ice is heated at a constant rate for $16$ minutes. For the first $2$ minutes, the temperature rises uniformly with time, then it remains constant for next $10$ minutes and again the temperature rises at uniform rate for last $4$ minutes. Final temperature at the end of $16$ minutes (in $^\circ$C), is :- [$C_{\mathrm{ice}}=0.5$ cal/g$^\circ$C, $C_{\mathrm{water}}=1$ cal/g$^\circ$C, $L_f=80$ cal/g]
\qn{2} A conducting rod of length $5$ m slides in a vertical plane against two insulated perpendicular surfaces as shown. Outward uniform magnetic field $B_0=3$ tesla exists in the region. If end $A$ of the rod is being pulled with constant speed $V_0=2\ \mathrm{m/s}\,\hat i$, then find the magnitude of induced Emf (in volt) across end $A$ and $B$ of the rod at the instant when the rod makes an angle of $\theta=37^\circ$ with horizontal.
\qfig{p1_phy_s3_q02_1}
\qn{3} A conductor $PQR$ of length $0.6$ metre and with a resistance $r=1\Omega$ moves on two parallel conducting rails $CD$ and $AB$ with a uniform speed $V=4$ m/s normal to a uniform magnetic field of induction $B=5$ T. The rails are connected by two resistance as shown in figure. If mechanical power needed for the motion is $3x$ watt, then find out the value of $x$.
\qfig{p1_phy_s3_q03_1}
\qn{4} A capacitor and a resistor are joined in parallel to an AC source as shown in the figure. The phase difference between currents $I$ and $I_1$ is $53^\circ$. Find the value of $\dfrac{1}{\omega C}$ (in SI unit).
\qfig{p1_phy_s3_q04_1}
\qn{5} A capacitor of capacitance $1\ \mu$F is charged to potential difference $2$ V and is connected to an inductor of $1$ mH. At an instant when potential difference across the capacitor is $1$ V, the current in the circuit is $10^{-2}\sqrt{\dfrac{10x}{3}}$ ampere. Find out the value of $x$.
\qn{6} Consider two blocks of mass $m$ and $2m$ which are connected by an ideal spring of spring constant $k$. Complete arrangement is placed on smooth horizontal plane. At $t=O$ velocities $3$ m/s and $6$ m/s are given to $m$ and $2m$ as shown in figure. If $t=t_1$ and $t=t_2$ is the first time instants ($t_1$ and $t_2\neq0$) when block of mass $m$ have maximum velocity and $2m$ have minimum velocity respectively then value of $|t_1-2t_2|$ comes out to be $\pi\sqrt{\dfrac{2m}{Nk}}$, find $N$.
\qfig{p1_phy_s3_q06_1}
\qn{7} Consider a ring of mass $M=1$ kg and radius $R$ which is connected to an inextensible string which passes over a frictionless pulley. The other end of string is connected to upper end of a massless spring having spring constant $2k$ ($k=25$ N/cm). Ring can rotate in the vertical plane about hinge without any friction. If horizontal position of ring is in equilibrium position then time period of small oscillation of the ring comes out $T=\dfrac{\pi}{100}\sqrt N$s, find value of $N$.
\qfig{p1_phy_s3_q07_1}
\qn{8} Two vibrating strings of the same material but length $L$ and $2L$ have radii $2r$ and $r$ respectively. They are stretched under the same tension. Both the strings vibrates in their fundamental modes, the one of length $L$ with frequency $f_1$ and the other with frequency $f_2$. Then the ratio $\left(\dfrac{f_1}{f_2}\right)$ is :-
\papersec{2}
\partsec{1}{PHYSICS}
\sect{SECTION-I(i)}
\paraq{1}{2}
People waiting near the tracks of an electrically driven train have sometimes noticed a tingling if they touch a conducting object, such as a pipe, connected to the ground. What causes the tingling? The reason is that the electrically driven train is connected to power station through overhead lines which carry Alternating Current. Because such current continuously changes direction and strength, the magnetic field it produces in surrounding space also changes in direction and strength. In our body, this produces eddy currents, but it is too small for a person to notice. However, if the person touches a large conducting object, such as a pipe, the current can be large and noticeable.\par
Let us try to understand this phenomena. Suppose that the man is standing vertically and his hand is touching the pipe. The overhead wire is an infinite wire carrying AC. Assume that man, pipe and the ground made a rectangular conducting loop of resistance $500\Omega$, dimensions as shown. The current in overhead wire is $100\sin400t$ (A).
\qfig{p2_phy_s1_para1_2_1}
\qn{1} In which of the situations, a current will flow through the loop ?
\begin{qoptions}
\opt{A}{\qfigi{p2_phy_s1_q01_A}}
\opt{B}{\qfigi{p2_phy_s1_q01_B}}
\opt{C}{\qfigi{p2_phy_s1_q01_C}}
\opt{D}{\qfigi{p2_phy_s1_q01_D}}
\end{qoptions}
\qn{2} Select the correct statement(s) for the induced current.
\begin{qoptions}
\opt{A}{The induced current will have the same phase as that of over head wire.}
\opt{B}{The induced current will have a maximum magnitude of $8\ell n2\ \mu$A.}
\opt{C}{The induced current will have phase difference of $90^\circ$ with over head wire.}
\opt{D}{The maximum magnitude of induced current is $4\ell n2\ \mu$A.}
\end{qoptions}
\paraq{3}{4}
A conducting slider $PQ$ of mass $m$ and length $\ell$ is free to translate on two vertical smooth rails in the presence of gravity. The slider is connected with an initially relaxed, non conducting spring of spring constant $k$ in uniform magnetic field $B$ into the plane. If the slider is released from rest and capacitor is connected in circuit then
\qfig{p2_phy_s1_para3_4_1}
\qn{3} Motion of rod will be
\begin{qoptions}
\opt{A}{SHM with angular frequency $\sqrt{\dfrac km}$}
\opt{B}{Oscillatory but not SHM}
\opt{C}{SHM with angular frequency $\sqrt{\dfrac{k}{m+B^2\ell^2C}}$}
\opt{D}{Amplitude of oscillation of rod will be $\dfrac{mg}{k}$}
\end{qoptions}
\qn{4} If suddenly spring is cut while rod is in motion, then :
\begin{qoptions}
\opt{A}{the rod will still perform oscillatory motion.}
\opt{B}{the rod will move downward and aquire terminal velocity.}
\opt{C}{acceleration of rod will be constant.}
\opt{D}{charge on capacitor will increase with time.}
\end{qoptions}
\paraq{5}{6}
Consider a square loop of wire with dimension $a$. One side of which is parallel to another long straight wire carrying current $i_0$, at a distance $a$. The resistance of the loop is $R$. Neglect the self inductance of the loop. Initially switch $S$ was closed and at $t=0$ it is opened.
\qfig{p2_phy_s1_para5_6_1}
\qn{5} What is the net amount of charge that will flow after $t=0$ through a fixed cross-section of the wire of the loop?
\begin{qoptions}
\opt{A}{$\dfrac{\mu_0i_0\ell n(2)\,a}{\pi R}$}
\opt{B}{$\dfrac{2\mu_0i_0\ell n(2)\,a}{\pi R}$}
\opt{C}{$\dfrac{\mu_0i_0\ell n(2)\,a}{2\pi R}$}
\opt{D}{$\dfrac{4\mu_0i_0\ell n(2)\,a}{\pi R}$}
\end{qoptions}
\qn{6} What is the net momentum $P$ given to the loop during the switching-off of the current :-
\begin{qoptions}
\opt{A}{$\dfrac{\mu_0^2i_0^2\times\ell n(2)\,a}{2\pi^2R}$}
\opt{B}{$\dfrac{\mu_0^2i_0^2\times\ell n(2)\,a}{4\pi^2R}$}
\opt{C}{$\dfrac{\mu_0^2i_0^2\times\ell n(2)\,a}{\pi^2R}$}
\opt{D}{None}
\end{qoptions}
\sect{SECTION-I(ii)}
\paraq{7}{8}
A simple pendulum of length $L=\dfrac52$m has a bob of mass $m=1$kg. The bob is connected to light horizontal spring of force constant $k=5$ N/m. The spring is relaxed when pendulum is vertical (see fig 1.) Now consider a similar situation in which spring is replaced with elastic cord of same force constant $k$. The cord is relaxed when the pendulum is vertical (see fig. 2). The bob is pulled slightly in same plane and released in both cases. (Use $\pi=3.14$, $g=10\ \mathrm{m/s^2}$)
\qfig{p2_phy_s2_para7_8_1}
\qn{7} Time period (in sec.) of small oscillations in Fig. 1\underline{\hspace{1.5cm}}
\qn{8} Time period (in sec.) of small oscillations in Fig. 2\underline{\hspace{1.5cm}}
\paraq{9}{10}
In the RLC series circuit shown, the readings of voltmeters are $V_1=150$ V, $V_2=50$ V and the source has emf $130$V.
\qfig{p2_phy_s2_para9_10_1}
\qn{9} The power factor of the circuit is $\dfrac{6n}{13}$ find $n$.
\qn{10} Value of $\dfrac{4V_L}{V_C}$ is:
\paraq{11}{12}
Two strings $AB$ and $BC$ of length $1$ meter and $1/2$ meter respectively are joined together at end $B$ as shown. Point $B$ is the junction of both strings. The mass per unit length of string $AB$ and $BC$ are $0.16$ kg/m and $0.25$ kg/m respectively. A tension of $100$ N is maintained in both the strings. A sinusoidal transverse wave pulse of amplitude $2$mm is produced at end `$A$'. It is partly reflected and partly transmitted at joint $B$.
\qfig{p2_phy_s2_para11_12_1}
\qn{11} Find value of $\dfrac{A_t}{A_i}$ (where $A_t\to$ Amplitude of transmitted wave, $A_i\to$ amplitude of incidence wave)
\qn{12} Calculate the value of $\dfrac{P_{\mathrm{transmitted}}}{P_{\mathrm{incident}}}$ (where $P_{\mathrm{transmitted}}\to$ power transmitted, $P_{\mathrm{incident}}\to$ power incident)
\sect{SECTION-II(i)}
\qn{1} A transverse periodic wave on a string with linear mass density of $0.2$ kg/m is described by equation.\\
$y=0.01\sin\{220t-11x\}$\\
where $x$ and $y$ are in metres and $t$ in seconds. The tension (in N) in the string is equal to :
\qn{2} A sinusoidal voltage of peak value $250$ V is applied to a series LCR circuit, in which $R=8\Omega$, $L=24$ mH and $C=60\mu$F. The value of power dissipated at resonant condition is `$x$' kW. The value of $x$ to the nearest integer is \underline{\hspace{1.5cm}}.
\qn{3} A linear harmonic oscillator is subjected to a periodic force such that it's amplitude is maximum at an angular frequency of $10$ rad/s. Phase difference between applied force and it's velocity at this frequency is $\phi$. It's damping coefficient is $0.1$ / sec. Fill the value of $|\cot\phi|$.
\sect{SECTION-IV}
\qn{1} A series LCR circuit is connected to a $45\sin(\omega t)$ Volt source. The resonant angular frequency of the circuit is $10^5\ \mathrm{rad\,s^{-1}}$ and current amplitude at resonance is $I_0$. When the angular frequency of the source is $\omega=8\times10^4\ \mathrm{rad\,s^{-1}}$, the current amplitude in the circuit is $0.05\,I_0$. If $L=50$ mH, match each entry in column-I with an appropriate value from column-II and choose the correct option.\par\smallskip
{\footnotesize\setlength{\tabcolsep}{2pt}\begin{tabular}{|l|p{3.7cm}|l|l|}
\hline
\multicolumn{2}{|c|}{Column-I} & \multicolumn{2}{c|}{Column-II}\\
\hline
(A) & $I_0$ in mA & (P) & $44.4$\\
\hline
(B) & The quality factor of the circuit & (Q) & $18$\\
\hline
(C) & The bandwidth of the circuit in $\mathrm{rad\,s^{-1}}$ & (R) & $400$\\
\hline
(D) & The peak power dissipated at resonance in Watt & (S) & $2250$\\
\hline
 & & (T) & $500$\\
\hline
\end{tabular}}
\qn{2} Here $I$ and $V$ represent current and voltage respectively (steady state for DC and rms for AC) match Column-I with Column-II using given codes.
\end{multicols}
\noindent\begin{minipage}{\textwidth}
\centering
\setlength{\tabcolsep}{4pt}
\begin{tabular}{|c|c|c|l|}
\hline
\multicolumn{2}{|c|}{Column-I} & \multicolumn{2}{c|}{Column-II}\\
\hline
(A) & \qfigi{p2_phy_s4_q02_A} & (P) & $I=0.2$ A\\
\hline
(B) & \qfigi{p2_phy_s4_q02_B} & (Q) & $I=0$ A\\
\hline
(C) & \qfigi{p2_phy_s4_q02_C} & (R) & $V=20$ Volt\\
\hline
(D) & \qfigi{p2_phy_s4_q02_D} & (S) & $V=0$ Volt\\
\hline
 & & (T) & $V=10$ Volt\\
\hline
\end{tabular}
\end{minipage}
\begin{multicols}{2}
\end{multicols}
\end{document}
